\begin{frame}{Conclusões}
    \framesubtitle{Síntese dos Achados Principais}
    \begin{block}{Hipótese Confirmada}
        Análise de redes complexas revela estruturas competitivas latentes não capturadas por métricas tradicionais de contagem de medalhas.
    \end{block}

    \vspace{0.3cm}
    \begin{columns}[T]
        \begin{column}{0.48\textwidth}
            \textbf{Diferenças Estruturais:}
            \begin{itemize}
                \footnotesize
                \item Esportes individuais: alta modularidade (0.73)
                \item Esportes coletivos: baixa modularidade (0.003)
                \item Comunidades refletem múltiplos fatores combinados
            \end{itemize}
        \end{column}
        \begin{column}{0.48\textwidth}
            \textbf{Métricas de Rede:}
            \begin{itemize}
                \footnotesize
                \item PageRank captura qualidade de competidores
                \item Atletas-ponte conectam comunidades
                \item Estrutura por gênero é diferenciada
            \end{itemize}
        \end{column}
    \end{columns}
\end{frame}

\begin{frame}{Contribuições do Trabalho}
    \framesubtitle{Avanços Metodológicos e Analíticos}
    \begin{enumerate}
        \item \textbf{Metodologia Validada}
        \begin{itemize}
            \footnotesize
            \item Construção de redes competitivas direcionadas ponderadas
            \item CDF dos pesos e PageRank confirmam premissas da modelagem
            \item 0 nós isolados garantidos através de filtragem
        \end{itemize}

        \item \textbf{Caracterização Multidimensional}
        \begin{itemize}
            \footnotesize
            \item 36 comunidades detectadas em 4.659 atletas
            \item Distribuição equilibrada de medalhas (31.8\%/33.4\%/34.8\%)
            \item Análise temporal via Entropia de Shannon
        \end{itemize}

        \item \textbf{Insights Estruturais}
        \begin{itemize}
            \footnotesize
            \item Comunidades emergem de fatores complexos
            \item Diversidade temporal varia por tipo de esporte
            \item Diferenças estruturais entre modalidades quantificadas
        \end{itemize}
    \end{enumerate}
\end{frame}

\begin{frame}{Limitações e Escolhas Metodológicas}
    \framesubtitle{Considerações para Interpretação}
    \begin{itemize}
        \item \textbf{Modelagem de Rivalidades Diretas}
        \begin{itemize}
            \footnotesize
            \item Agrupamento (Ano, Evento) captura \textbf{competições reais} no mesmo pódio
            \item Comparação cross-geracional (ex: Phelps vs Spitz) requer modelagem alternativa
            \item Trade-off: precisão histórica vs. riqueza relacional
        \end{itemize}

        \item \textbf{Amostra de Esportes}
        \begin{itemize}
            \footnotesize
            \item Análise expandida para 6 esportes (Swimming, Basketball, Football, Athletics, Judo, Boxing)
            \item Cobertura de tipos estruturais: individual, coletivo e misto
        \end{itemize}

        \item \textbf{Dados até 2021}
        \begin{itemize}
            \footnotesize
            \item Dataset exclui Paris 2024
            \item 125 anos de cobertura olímpica
        \end{itemize}
    \end{itemize}
\end{frame}

\begin{frame}{Trabalhos Futuros}
    \framesubtitle{Direções de Pesquisa}
    \begin{enumerate}
        \item \textbf{Expansão para Mais Modalidades}
        \begin{itemize}
            \footnotesize
            \item Esportes de combate, eventos cronometrados, modalidades artísticas
            \item Validação de propriedades estruturais gerais
        \end{itemize}

        \item \textbf{Análise de Evolução Topológica}
        \begin{itemize}
            \footnotesize
            \item Snapshots temporais de redes por década
            \item Estabilidade de comunidades e persistência de hierarquias
        \end{itemize}

        \item \textbf{Comparação com Rankings Oficiais}
        \begin{itemize}
            \footnotesize
            \item Validação de PageRank contra sistemas de classificação estabelecidos
            \item Identificação de discrepâncias estruturais
        \end{itemize}

        \item \textbf{Análise de Gênero Sistematizada}
        \begin{itemize}
            \footnotesize
            \item Comparação sistemática de propriedades estruturais
            \item Padrões de equidade e disparidade ao longo da história
        \end{itemize}
    \end{enumerate}
\end{frame}

\begin{frame}{Possibilidades de Modelagem Alternativa}
    \framesubtitle{Explorando Múltiplas Perspectivas}
    \begin{block}{Redes Bipartidas (Atletas $\leftrightarrow$ Eventos)}
        \footnotesize
        \begin{itemize}
            \item Atletas conectados a eventos onde conquistaram medalhas
            \item Projeção em atletas: rivalidade indireta (competiram nos mesmos eventos)
            \item Métrica de versatilidade: atletas que dominam múltiplos eventos
        \end{itemize}
    \end{block}

    \vspace{0.2cm}
    \begin{block}{Redes Multiplex (Múltiplas Camadas)}
        \footnotesize
        \begin{itemize}
            \item Camada 1: Competição direta (Year+Event) -- rivalidades reais
            \item Camada 2: Domínio histórico (Event apenas) -- comparação cross-geracional
            \item Camada 3: Mesma nacionalidade -- estrutura de times olímpicos
            \item Análise inter-camadas: múltiplas dimensões de rivalidade
        \end{itemize}
    \end{block}

    \vspace{0.2cm}
    \begin{alertblock}{Métricas Alternativas}
        \footnotesize
        \centering
        Hub \& Authority (HITS), Dominance Score, Temporal PageRank
    \end{alertblock}
\end{frame}

\begin{frame}{Mensagem Final}
    \framesubtitle{Aplicação de Redes Complexas ao Esporte Olímpico}
    \begin{center}
        \Large
        \textbf{Análise de Redes Complexas revela estruturas competitivas latentes e padrões evolutivos invisíveis em abordagens tradicionais.}
    \end{center}

    \vspace{0.5cm}
    \begin{block}{}
        \centering
        Este trabalho estabelece fundamentos metodológicos sólidos para aplicação de ciência de redes ao contexto esportivo, com potencial para avanço teórico e aplicações práticas em gestão e desenvolvimento atlético.
    \end{block}

    \vspace{0.5cm}
    \begin{center}
        \Large
        \textbf{Obrigado!}
    \end{center}
\end{frame}
